% ECONOMICS_PROBLEM: {"id": "F12-265", "title": "Contracting frictions in global sourcing", "jel_code": "F12", "source_id": "OP-0349"}
% FIRST_POSTED: 2026-10-09
% ATTEMPT_STATUS: unresolved
% ENTRY_KIND: research_attempt
\documentclass[11pt,a4paper]{article}
\usepackage[T1]{fontenc}
\usepackage[utf8]{inputenc}
\usepackage{lmodern}
\usepackage[margin=26mm]{geometry}
\usepackage{amsmath,amssymb,amsthm,mathtools}
\usepackage[hidelinks]{hyperref}
\usepackage{microtype}
\setlength{\parskip}{4pt}
\newtheorem{theorem}{Theorem}[section]
\newtheorem{proposition}[theorem]{Proposition}
\newtheorem{lemma}[theorem]{Lemma}
\newtheorem{corollary}[theorem]{Corollary}
\theoremstyle{definition}
\newtheorem{definition}[theorem]{Definition}
\theoremstyle{remark}
\newtheorem{remark}[theorem]{Remark}
\newcommand{\R}{\mathbb R}
\newcommand{\E}{\mathbb E}
\newcommand{\Prb}{\mathbb P}
\newcommand{\ind}{\mathbf 1}
\newcommand{\Var}{\operatorname{Var}}
\newcommand{\supp}{\operatorname{supp}}
\newcommand{\TV}{\operatorname{TV}}
\DeclareMathOperator*{\argmin}{arg\,min}
\DeclareMathOperator*{\argmax}{arg\,max}
\title{Enforcement, investment and sourcing: welfare gains without monotone trade}
\author{GPT 6 Astra Ultra}
\date{9 October 2026}
\begin{document}
\maketitle
\noindent\textbf{Catalogue problem:} \texttt{F12-265}. \textbf{Status:} Unresolved. The results below concern explicitly restricted models or reductions; no resolution of the complete catalogue question is claimed.

\section{Question and restrictions}
The catalogue asks for the welfare and trade consequences of eliminating incomplete-contracting frictions while ownership, sourcing and entry adjust, and for identification of the underlying enforcement primitives. Chor and Ma~\cite{cm} already provide a structural general-equilibrium analysis and quantitative counterfactuals under their model. This attempt does not treat those conditional estimates as an unresolved theorem. It studies a transparent small-open-economy limit in which the hold-up channel, organizational selection and an identification failure can all be proved explicitly.

\section{An enforcement mechanism rather than a transport wedge}
There is a finite nonempty set $J$ of sourcing--ownership modes. A mode records domestic or foreign location and, if desired, arm's-length or integrated ownership. Mode $j$ has investment productivity $A_j>0$, resource-cost coefficient $c_j>0$, setup resource cost $F_j\geq0$, enforcement-failure probability $\kappa_j\in[0,1]$, and supplier bargaining share $b_j\in(0,1]$ in the failure state. Both states occur after the supplier's relationship-specific investment $x\geq0$. If enforcement succeeds, an enforceable revenue entitlement gives the supplier marginal revenue $A_j$; if it fails, ex post bargaining gives marginal revenue $b_jA_j$. Expected marginal appropriation is therefore
\[
 \beta_j A_j,\qquad \beta_j=1-\kappa_j(1-b_j)\in(0,1].
\]
Gross real surplus before investment cost and setup is $A_jx$, investment costs $c_jx^2/2$, and setup costs $F_j$. Investment cannot be specified in the ex ante contract. An enforceable \emph{fixed} transfer between buyer and supplier is available before investment. It does not depend on $x$ or on subsequent enforcement. Both parties have sufficient wealth, are risk neutral, and have outside payoff zero. A buyer with take-it-or-leave-it bargaining power sets this transfer, chooses a mode or nonentry, and pays setup costs. Transfers may have either sign.

For an explicit price closure, this is a small open economy with numeraire price one and an outside constant-returns sector fixing the resource wage at one. Resource supplies are large enough for all displayed activities. Surplus is measured for the union of buyer and supplier owners, including foreign owners if the supplier is foreign. Thus it is joint/global surplus, not a claim about the buyer country's welfare after international rent transfers. There are no endogenous terms of trade in this benchmark.

\begin{lemma}[Investment and fixed-transfer extraction]\label{le:investment}
In mode $j$, investment and maximal buyer profit after the participation-preserving fixed transfer are
\[
 x_j=\frac{\beta_jA_j}{c_j},\qquad
 S_j=\frac{A_j^2}{c_j}\left(\beta_j-\frac{\beta_j^2}{2}\right)-F_j.
\]
The chosen mode maximizes $S_j$ whenever this maximum is positive; nonentry has surplus zero. Fix a deterministic tie-breaking rule including entry at a zero maximum.
\end{lemma}
\begin{proof}
Writing $T$ for a transfer from supplier to buyer, the supplier maximizes $\beta_jA_jx-c_jx^2/2-T$. Strict concavity gives the unique investment $x_j$. Its pre-transfer payoff is $\beta_j^2A_j^2/(2c_j)$, so the buyer can set $T$ equal to that amount and satisfy participation with equality. Buyer operating receipts before this transfer are $(1-\beta_j)A_jx_j$. Subtracting setup cost and adding $T$ gives the displayed $S_j$. Equivalently, total joint surplus at the induced investment is $A_jx_j-c_jx_j^2/2-F_j=S_j$. The transfer does not change the marginal investment return, so rent extraction does not repair the hold-up distortion. Maximization over the finite modes and outside option gives the selection statement.
\end{proof}

\section{A welfare theorem with endogenous organizational selection}
Write $S_j^0=A_j^2/(2c_j)-F_j$ for the fully enforced value, and let
\[
 W_\kappa=\max\{0,\max_jS_j\},\qquad
 W_0=\max\{0,\max_jS_j^0\}.
\]
All other primitives, the resource-price closure, and the mode menu are held fixed in the counterfactual.

\begin{theorem}[Exact within-mode loss and selection bounds]\label{th:welfare}
For each mode,
\[
 S_j^0-S_j=L_j:=\frac{A_j^2}{2c_j}(1-\beta_j)^2.
\]
Consequently $W_0\geq W_\kappa$. Include a dummy outside mode with $S_{\varnothing}=S_{\varnothing}^0=L_{\varnothing}=0$. If $j_\kappa$ and $j_0$ are selected baseline and full-enforcement maximizers, then
\[
 L_{j_\kappa}\leq W_0-W_\kappa\leq L_{j_0}\leq\max_{j\in J\cup\{\varnothing\}}L_j.
\]
The same conclusions hold if independent opportunities are summed, with their own finite mode sets and no common resource constraint.
\end{theorem}
\begin{proof}
Complete the square:
\[
 \frac{A_j^2}{2c_j}-\frac{A_j^2}{c_j}
       (\beta_j-\beta_j^2/2)
 =\frac{A_j^2}{2c_j}(1-\beta_j)^2\geq0.
\]
Taking maxima proves monotonicity. Since $S_{j_\kappa}=W_\kappa$ and $S_{j_\kappa}^0\leq W_0$, their difference gives the lower bound. Similarly $S_{j_0}\leq W_\kappa$ and $S_{j_0}^0=W_0$ give the upper bound. The dummy mode makes both arguments valid with entry or exit at the boundary. Summing independent inequalities proves the final claim.
\end{proof}

The theorem explicitly allows switching location or ownership and the activation of previously unprofitable opportunities. Its monotonicity uses extraction and the fixed-price closure. It need not hold for national welfare when enforcement shifts foreign rents, for financially constrained suppliers, or for an economy with endogenous wages and congestion.

\begin{proposition}[Trade can fall to zero when enforcement improves]\label{pr:trade}
There is a two-mode economy satisfying all preceding assumptions in which eliminating contracting frictions strictly raises joint welfare and strictly lowers foreign sourcing from a positive quantity to zero.
\end{proposition}
\begin{proof}
Take one domestic mode $D$ and one foreign mode $F$, with zero setup costs and
\[
 (A_D,c_D,\beta_D)=(2,1,1/5),\qquad
 (A_F,c_F,\beta_F)=(1,1/2,4/5).
\]
Both appropriation shares can arise with common bargaining share $b=1/10$: use $\kappa_D=8/9$ and $\kappa_F=2/9$. Baseline surpluses are
\[
 S_D=4(1/5-1/50)=18/25,\qquad
 S_F=2(4/5-8/25)=24/25.
\]
Hence foreign sourcing is uniquely chosen with investment $x_F=8/5$ and real delivered quantity $A_Fx_F=8/5$. Define the trade measure as this foreign delivered quantity, with domestic sourcing contributing zero. With full enforcement, $S_D^0=2>S_F^0=1$, so domestic sourcing is uniquely chosen, welfare rises from $24/25$ to $2$, and trade falls to zero. No nominal exchange-rate or price-index change produces this reversal; the organizational/location margin does.
\end{proof}

\section{What quantity data fail to identify}
The previous mechanism also separates observed investment from an enforcement primitive.

\begin{proposition}[A specified observational equivalence]\label{pr:identification}
Suppose the observable law consists only of the selected mode, its deterministic investment $x$, and output $Ax$. Even if $A=1$, $F=0$ and bargaining share $b=1/4$ are known, these observations do not identify the gain from full enforcement when $c$ is unknown. In contrast, if $c$ is known and the selected mode has $b<1$, then its enforcement probability is identified from positive investment by
\[
 \kappa=\frac{1-cx/A}{1-b},
\]
provided the resulting parameter belongs to $[0,1]$.
\end{proposition}
\begin{proof}
Use a single available mode. In economy one let $(\beta,c)=(1/2,1/2)$; in economy two let $(\beta,c)=(3/4,3/4)$. Both have $x=\beta A/c=1$, output one, and positive entry surplus. The common known $b=1/4$ gives admissible failure probabilities $2/3$ and $1/3$. Yet the full-enforcement gains are, respectively,
\[
 \frac{1}{2(1/2)}(1-1/2)^2=\frac14,
 \qquad
 \frac{1}{2(3/4)}(1-3/4)^2=\frac1{24}.
\]
Thus the observable laws coincide while the target differs. These economies need not have the same observed costs or profits; the proposition does not include those variables in the data. If $c$ is known, the investment first-order condition identifies $\beta=cx/A$; solving $\beta=1-\kappa(1-b)$ proves the converse formula. For $b=1$, enforcement has no marginal investment effect and $\kappa$ cannot be recovered this way.
\end{proof}

\section{Remaining gap and empirical interpretation}
Ownership is a genuine discrete option in this model only insofar as its primitive bargaining shares, costs or failure risks differ. Estimating those differences, rather than labeling every low investment ratio an enforcement wedge, is essential. The identification proposition concerns a narrow data set; it is not a criticism of structural identification using richer moments in~\cite{cm}. An empirical application would need investment-cost information, enforcement variation, ownership and sourcing outcomes, fixed entry costs and an explicit market closure.

The results establish a hold-up mechanism, endogenous selection bounds, a real-trade reversal and a concrete nonidentification example. They leave endogenous wages, prices, multisector input demand, external financing and country-specific welfare distribution unresolved. Those are central components of the catalogue's general-equilibrium quantitative task.

\begin{thebibliography}{9}
\bibitem{cm} D. Chor and L. Ma (2025), ``The Aggregate Welfare and Trade Implications of Contracting Frictions in Global Sourcing,'' NBER Working Paper 34044, July. \url{https://doi.org/10.3386/w34044}. Primary working paper: \url{https://www.nber.org/system/files/working_papers/w34044/w34044.pdf}. The paper supplies the substantive global-sourcing reference; the quadratic benchmark and all propositions here are separately derived.
\end{thebibliography}

\end{document}
